\section{构造原理与进一步的问题}

各项点数来自有限结构的具体改变。在码构造中，交换收益由冲突并集决定；
允许局部符号码替换，则进一步扩大了可用的变化。在赤道构造中，转向另一格壳
提供了与全部既有帽层相容的新方向。在截面构造中，母格中的嵌入决定了投影重数，
而这些信息仅从抽象 Gram 型看不出来。

这些认识也指出了进一步研究的具体变量。支持族可以与符号分配共同改变；
头方向类和尾标签可以按其对完整构型的净贡献共同选择；
嵌入截面可以借助更大的母根系研究；
矩恒等式也可以针对最容易给出见证的统计量选择，
而不必总是重建每一个纤维的全部重数。

配套有限数据将这些选择保存为明确的数学对象。
读者可以据此重新考察构造、测试新的局部变化并计算其对终值的影响，
无需重复整个发现搜索。
